\psset{unit=1cm,arrowsize=2pt 3}
\begin{center}
\begin{pspicture}(-4,-2)(4,2.3)

\psline(-4,-2)\psline(3.9,-0.7)\psline(0,2.3)
\psline[linecolor=blue]{->}(-1.1,-0.55)\psline[linecolor=blue]{->}(1.7,-0.285)\psline[linecolor=blue]{->}(0,1.8)
\pspolygon(-3.5,-1.73)(3.4,-0.6)(0,1.8)
\psline(0,1.8)(1.55,-0.9)
\uput[ul](-3.5,-1.73){A} \uput[ur](3.4,-0.6){B} \uput[ul](0,1.8){C} \uput[d](0,0){O} \uput[dr](1.55,-0.9){K}
\uput[u](-0.5,-0.3){\blue $\vect{\imath}$} \uput[u](0.6,-0.1){\blue $\vect{\jmath}$} \uput[l](0,0.8){\blue $\vect{k}$}
\end{pspicture}
\end{center}

\medskip

L'espace est muni d'un repère orthonormé \Oijk.

On considère les points
A\,$\left(2\sqrt 3~;~0~;~0\right)$, B\,(0~;~2~;~0), C\,(0~;~0~;~1) et  K\,$\left(\frac{\sqrt 3}{2}~;~\frac32~;~0\right)$.




\begin{enumerate}
\item 
La droite (CK) est l'ensemble des points M\,$\left (x\;,\; y \;,\; z\right )$ tels que $\vectt{CM}$ et $\vectt{CK}$ soient colinéaires, c'est-à-dire tels que $\vectt{CM} = t.\vectt{CK}$ où $t\in\R$.

$\vectt{CM}$ a pour coordonnées  
$\begin{pmatrix}
x_{\text{M}} - x_{\text{C}} \\ y_{\text{M}} - y_{\text{C}} \\ z_{\text{M}} - z_{\text{C}}
\end{pmatrix}
=
\begin{pmatrix}
x\\ y \\ z -1
\end{pmatrix}$
et celles de $\vectt{CK}$ sont
$\begin{pmatrix}
x_{\text{K}} - x_{\text{C}} \\ y_{\text{K}} - y_{\text{C}} \\ z_{\text{K}} - z_{\text{C}}
\end{pmatrix}
=
\begin{pmatrix}
\frac{\sqrt{3}}{2}\\[5pt] \frac{3}{2} \\[5pt] -1
\end{pmatrix}$

$\vectt{CM} = t.\vectt{CK}
\iff
\begin{pmatrix}
x\\ y \\ z-1
\end{pmatrix}
=
t.\begin{pmatrix}
\frac{\sqrt{3}}{2}\\[5pt] \frac{3}{2} \\[5pt] -1
\end{pmatrix}
\iff
\left \{
\begin{array}{l !{=} l}
x & \frac{\sqrt{3}}{2} t\\[5pt]
y & \frac{3}{2} t\\[5pt]
z & -t+1
\end{array}
\right .
$

La droite (CK) a donc pour représentation paramétrique 
$\left\{\begin{array}{l !{=} l}
x& \frac{\sqrt 3}{2}t\\[5pt]
y& \frac32 t\\[5pt]
y& - t + 1
\end{array}\right. \: (t \in \R)$

\item Soit M$(t)$ un point de la droite (CK) paramétrée par un réel $t$.
Le point M a donc pour coordonnées $\left (\frac{\sqrt{3}}{2}t\;,\; \frac{3}{2}t\;,\; -t+1\right )$

$\begin{aligned}
\text{OM}(t)^2
&  = \left ( x_{\text{M}} - x_{\text{O}}\right )^2 + \left ( y_{\text{M}} - y_{\text{O}}\right )^2 + \left ( z_{\text{M}} - z_{\text{O}}\right )^2 
=\left ( \dfrac{\sqrt{3}}{2}t\right )^2 + \left (\dfrac{3}{2} t\right )^2 + \left (-t+1\right )^2\\
& = \dfrac{3}{4}t^2 +\dfrac{9}{4}t^2 +t^2-2t+1
= 4t^2 -2t +1
\end{aligned}$

Donc: OM$(t) = \sqrt{4t^2-2t+1}$.



\item Soit $f$ la fonction définie et dérivable sur $\R$ par $f(t) = \text{OM}(t)$.
	\begin{enumerate}
		\item 
$f(t)=\sqrt{4t^2-2t+1}$ donc $f'(t)=\dfrac{8t-2}{2\sqrt{4t^2-2t+1}}$		

On étudie le signe de $f'(t)$ sur $\R$.
		
\[ \begin{tablvar}[intervalwidth=6em]{2}
 \hline
 t & -\infty & & \frac{1}{4} & & +\infty\\
 \hline
8t-2 & & - & \barre[0] & + & \\
 \hline
2\sqrt{4t^2-2t+1} & & + & \barre & + & \\
 \hline
f'(t) & & - & \barre[0] & + & \\
 \hline
 \end{tablvar}\]
		
Donc la fonction $f$ est strictement décroissante sur l'intervalle  $\left ] -\infty\;,\; \frac{1}{4}\right ]$, et strictement croissante sur l'intervalle $\left [ \frac{1}{4}\;,\; +\infty \right [$.		
		
		\item 
On en déduit que $f$ atteint son minimum pour $t=\frac{1}{4}$.		
		
	\end{enumerate}
	
\item 
Le projeté orthogonal du point O sur la droite (CK) est le point M de la droite (CK) tel que la distance OM soit minimale; le point de la droite (CK) réalisant ce minimum correspond donc à $t=\frac{1}{4}$.

C'est donc le point de coordonnées $\left (\frac{\sqrt{3}}{2}\times \frac{1}{4}\;,\; \frac{3}{2}\times \frac{1}{4}\;,\; -\frac{1}{4}+1\right )$ soit $\left (\frac{\sqrt{3}}{8}\;,\; \frac{3}{8}\;,\; \frac{3}{4}\right )$ est le projeté orthogonal du point O sur la droite (CK); on l'appelle H.

\item
\begin{list}{\textbullet}{On montre que H appartient au plan (ABC).}
\item 
$\vectt{AK}$ a pour coordonnées  
$\begin{pmatrix}
\frac{\sqrt{3}}{2} - 2\sqrt{3}\\[5pt] \frac{3}{2}-0 \\[5pt] 0-0
\end{pmatrix}
=
\begin{pmatrix}
-\frac{3}{2}\sqrt{3} \\[5pt] \frac{3}{2} \\[5pt] 0-0
\end{pmatrix}$
et celles de $\vectt{AB}$ sont
$\begin{pmatrix}
0-2\sqrt{3} \\ 2-0 \\ 0-0
\end{pmatrix}
=
\begin{pmatrix}
-2\sqrt{3} \\ 2\\ 0
\end{pmatrix}$

Donc $\vectt{AK} = \frac{3}{4}\vectt{AB}$ donc les points A, K et B sont alignés donc $\text{K}\in \left (\text{AB}\right )$.

\item On a:
$\left .
\begin{array}{r}
\text{K} \in \text{(AB)}\\
\text{(AB)} \subset \text{(ABC)}
\end{array}
\right \}
\implies \text{K} \in  \text{(ABC)}
\implies \text{(CK)} \subset  \text{(ABC)}$

\item On a:
$\left .
\begin{array}{r}
\text{H} \in \text{(CK)}\\
\text{(CK)} \subset  \text{(ABC)}
\end{array}
\right \}
\implies \text{H} \in  \text{(ABC)}$
\end{list}

Donc H est un point du plan (ABC).

\begin{list}{\textbullet}{On montre que H est l'orthocentre du triangle ABC.}
\item 
$\vectt{CK}\cdot \vectt{AB} 
= \frac{\sqrt{3}}{8}\times \left ( -2\sqrt{3}\right ) + \frac{3}{8}\times 2 + \left (-\frac{1}{4} \right ) \times 0
= -\frac{3}{4}+\frac{3}{4}+0=0$
donc $\vectt{CK}\perp \vectt{AB}$ donc $\text{(CK)} \perp \text{(AB)}$.

Comme $\text{H} \in \text{(CK)}$, on en déduit que $\text{(CH)} \perp \text{(AB)}$.
\item 
$\vectt{AH}$ a pour coordonnées  
$\begin{pmatrix}
\frac{\sqrt{3}}{8} - 2\sqrt{3}\\[5pt] \frac{3}{8}-0 \\[5pt] \frac{3}{4}-0
\end{pmatrix}
=
\begin{pmatrix}
-\frac{15}{8}\sqrt{3} \\[5pt] \frac{3}{8} \\[5pt] \frac{3}{4}
\end{pmatrix}$
et celles de $\vectt{BC}$ sont
$\begin{pmatrix}
0-0 \\0- 2 \\ 1-0
\end{pmatrix}
=
\begin{pmatrix}
0 \\ -2\\ 1
\end{pmatrix}$

$\vectt{AH}\cdot \vectt{BC}
=  -\frac{15}{8}\sqrt{3} \times 0 +  \frac{3}{8} \times (-2)+  \frac{3}{4}\times 1
= 0-\frac{3}{4}+ \frac{3}{4} =0$
donc $\vectt{AH}\perp \vectt{BC}$ donc $\text{(AH)} \perp \text{(BC)}$.
\end{list}

(AH) et (CH) sont donc deux hauteurs du triangle (ABC) donc H est l'orthocentre de ce triangle.
\item 
	\begin{enumerate}
		\item 
\begin{list}{\textbullet}{On démontre que la droite (OH) est orthogonale au plan (ABC).}
\item On sait que $\vectt{OH} \perp \vectt{CK}$.
\item $\vectt{OH} \cdot \vectt{AB}
= \dfrac{\sqrt{3}}{8}\times \left (-2\sqrt{3}\right ) + \dfrac{3}{8}\times 2 + \dfrac{1}{4}\times 0
= -\dfrac{3}{4}+\dfrac{3}{4} + 0 = 0$ donc $\vectt{OH} \perp \vectt{AB}$.
\item Les vecteurs $\vectt{CK}$ et $\vectt{AB}$ ne sont pas colinéaires donc ce sont deux vecteurs directeurs du plan (ABC).
\end{list}		

Le vecteur $\vectt{OH}$ est orthogonal à deux vecteurs directeurs du plan (ABC) donc c'est un vecteur normal au plan (ABC). Donc la droite (OH) est orthogonale au plan (ABC).
		
		\item
D'après le cours, si le vecteur $\vect{v}$ de coordonnées $(a\,,\, b \,,\, c)$ est normal à un plan $\mathcal{P}$, alors ce plan a une équation cartésienne de la forme $ax+by+cz+d=0$.

Le vecteur $\vectt{OH}\;\left ( \frac{\sqrt{3}}{8}\,,\,\frac{3}{8}\,,\, \frac{3}{4}\right )$ est un vecteur normal au plan (ABC), donc le plan (ABC) a une équation de la forme $ \frac{\sqrt{3}}{8}x + \frac{3}{8}y + \frac{3}{4}z +d=0$ c'est-à-dire $x\sqrt{3} + 3y +6z +8d =0$.

On détermine la valeur de $d$ en exprimant que le point A appartient au plan (ABC).

$\begin{aligned}
\text{A}\in \text{(ABC)}
& \iff
x_{\text{A}}\sqrt{3}+3 y_{\text{A}} +6z_{\text{A}}+8d=0
\iff
2\sqrt{3}\times \sqrt{3} +3\times 0+6\times 0 + 8d=0\\
& \iff
6+8d=0
\iff 8d=-6
\end{aligned}$

Le plan (ABC) a pour équation: $x\sqrt{3} + 3y +6z -6 =0$.
	\end{enumerate}
	\item 
(CK) est une hauteur du triangle (ABC) donc l'aire de ce triangle est égale à $\frac{1}{2}\,\text{CK}\times \text{AB}$.

\begin{list}{\textbullet}{}
\item Le vecteur $\vectt{CK}$ a pour coordonnées $\left ( \frac{\sqrt{3}}{2}\,,\, \frac{3}{2}\,,\, -1\right )$ donc:

$\vectt{CK}^2
=\left ( \frac{\sqrt{3}}{2}\right )^2 + \left ( \frac{3}{2}\right )^2+\left ( -1\right )^2
= \frac{3}{4} + \frac{9}{4} +1
= 4$ donc $\text{CK}=2$.
\item Le vecteur $\vectt{AB}$ a pour coordonnées $\left ( -2\sqrt{3}\,,\, 2\,,\, 0\right )$ donc:

$\vectt{AB}^2
=\left ( -2\sqrt{3}\right )^2 + \left ( 2\right )^2+\left (0\right )^2
= 12+4+0
= 16$ donc $\text{AB}=4$.
\end{list}
	
L'aire du triangle (ABC) est donc égale, en unités d'aire, à: $\frac{1}{2}\times 2\times 4$ soit 4.
\end{enumerate}
